\begin{comment}
  Homework solution of Calculus 2 - Chapter 4 Integral Application
  Licensed under MIT license
\end{comment}

\documentclass[12pt]{article}
\usepackage[utf8]{inputenc}
\usepackage{amsmath}
\usepackage{geometry}

\geometry{a4paper, margin=1in}

\title{No. 8 - Aplikasi Integral}
\author{Ahmad Fakhrul Bawani}
\date{}

\begin{document}

\maketitle

\section*{Solusi Volume Benda Putar}

Diketahui daerah yang dibatasi oleh kurva $y = 5 \sin x \cos x$ pada interval $[0, \pi/2]$ diputar mengelilingi sumbu-$x$. Kita akan menghitung volumenya menggunakan metode cakram (\textit{disk method}).

\subsection*{1. Menyederhanakan Fungsi}
Untuk mempermudah integrasi, kita gunakan identitas trigonometri $2 \sin x \cos x = \sin(2x)$:
\[ y = 5 \sin x \cos x = \frac{5}{2} (2 \sin x \cos x) = \frac{5}{2} \sin(2x) \]

\subsection*{2. Rumus Volume Metode Cakram}
Karena sumbu putar adalah sumbu-$x$, maka volumenya adalah:
\[ V = \int_{a}^{b} \pi [f(x)]^2 \, dx \]
Substitusi fungsi dan batas interval $[0, \pi/2]$:
\[ V = \pi \int_{0}^{\pi/2} \left( \frac{5}{2} \sin(2x) \right)^2 \, dx = \pi \int_{0}^{\pi/2} \frac{25}{4} \sin^2(2x) \, dx \]

\subsection*{3. Menghitung Integral}
Gunakan identitas $\sin^2 \theta = \frac{1 - \cos(2\theta)}{2}$. Dalam hal ini $\theta = 2x$, sehingga $\sin^2(2x) = \frac{1 - \cos(4x)}{2}$:
\begin{align*}
	V &= \frac{25\pi}{4} \int_{0}^{\pi/2} \frac{1 - \cos(4x)}{2} \, dx \\
	V &= \frac{25\pi}{8} \int_{0}^{\pi/2} (1 - \cos(4x)) \, dx \\
	V &= \frac{25\pi}{8} \left[ x - \frac{1}{4} \sin(4x) \right]_{0}^{\pi/2}
\end{align*}

Evaluasi batas integrasi:
\begin{itemize}
  \item Untuk $x = \pi/2$: $(\pi/2) - \frac{1}{4} \sin(4 \cdot \frac{\pi}{2}) = \pi/2 - \frac{1}{4} \sin(2\pi) = \pi/2 - 0 = \pi/2$
  \item Untuk $x = 0$: $0 - \frac{1}{4} \sin(0) = 0$
\end{itemize}

\[ V = \frac{25\pi}{8} \left( \frac{\pi}{2} \right) = \frac{25\pi^2}{16} \]

\subsection*{Kesimpulan}
Jadi, volume benda putar yang dihasilkan adalah $\mathbf{\frac{25\pi^2}{16}}$ \textbf{satuan kubik}.

\end{document}